\subsection{Síntesis: Confrontación entre Objetivos y Resultados}\label{subsec:sintesis-resultados}

La \autoref{tab:sintesis_resultados} consolida la confrontación entre los resultados esperados declarados en el alcance del proyecto y los resultados efectivamente obtenidos, indicando la evidencia que respalda cada uno. La columna izquierda reproduce los resultados esperados de la \autoref{tab:objetivos_resultados}, de modo que el contraste sea directo.

% Resultados obtenidos por objetivo (columna derecha de la tabla de síntesis)
\newcommand\objetivoEspecificoAObtenido{
  \begin{itemize}[leftmargin=*, topsep=0pt, itemsep=0pt]
    \item Patrón BFF introducido en \texttt{paralegales-ui}; lógica crítica reubicada al servidor y no inspeccionable desde el cliente (\autoref{fig:resultado-bff-red}).
    \item Integración con la API de la Rama Judicial encapsulada como \textit{proxy} controlado en el BFF.
  \end{itemize}
}

\newcommand\objetivoEspecificoBObtenido{
  \begin{itemize}[leftmargin=*, topsep=0pt, itemsep=0pt]
    \item \textit{Frontend} organizado en nueve \textit{bounded contexts} (Nuxt Layers) con tema visual aislado (\autoref{fig:resultado-feature-tree}).
    \item \textit{Backend} con patrón \textit{Handler--Service--Repository} por interfaz (\autoref{fig:api-gateway}); validación de desacoplamiento detallada en el capítulo de pruebas.
  \end{itemize}
}

\newcommand\objetivoEspecificoCObtenido{
  \begin{itemize}[leftmargin=*, topsep=0pt, itemsep=0pt]
    \item Telemetría JSON con identificadores de correlación en \textit{frontend} y \textit{backend} (\autoref{fig:resultado-traza-correlacionada}).
    \item Dashboards, alarmas y canal de Telegram operativos; infraestructura como código en CloudFormation (\autoref{fig:dashboard-observabilidad}, \autoref{fig:dashboard-observabilidad-alerts}, \autoref{fig:telegram-alerts-observability}).
  \end{itemize}
}

\newcommand\objetivoEspecificoDObtenido{
  \begin{itemize}[leftmargin=*, topsep=0pt, itemsep=0pt]
    \item \textit{Hooks} de Lefthook operativos: \textit{pre-commit}, \textit{commit-msg} y \textit{post-merge} (\autoref{fig:resultado-commit-rechazado}).
    \item Pipeline de CI con validación diferencial en GitHub Actions (\autoref{fig:resultado-ci}).
    \item Suite de pruebas de reglas de seguridad de Firestore: detallada en el capítulo de pruebas.
  \end{itemize}
}

\newcommand\sintesisResultadosCaption{Confrontación entre los resultados esperados y los resultados obtenidos por objetivo específico. \hspace{1em}}

\begingroup
\renewcommand{\arraystretch}{1.2}
\begin{longtable}{|p{8cm}|p{8cm}|}
  \hline
  % Table Header
  \grayTableHeaderCell{8cm}{Resultado esperado} &
  \grayTableHeaderCell{8cm}{Resultado obtenido y evidencia} \\
  \hline

  \endfirsthead

  \hline
  \grayTableHeaderCell{8cm}{Resultado esperado} &
  \grayTableHeaderCell{8cm}{Resultado obtenido y evidencia} \\
  \hline
  \endhead

  \continuacionTablaFooter{2}
  \endfoot
  \endlastfoot

  % Table Body
  \tableCell{\objetivoEspecificoAResulados} &
  \tableCell{\objetivoEspecificoAObtenido} \\
  \hline

  \tableCell{\objetivoEspecificoBResultados} &
  \tableCell{\objetivoEspecificoBObtenido} \\
  \hline

  \tableCell{\objetivoEspecificoCResultados} &
  \tableCell{\objetivoEspecificoCObtenido} \\
  \hline

  \tableCell{\objetivoEspecificoDResultados} &
  \tableCell{\objetivoEspecificoDObtenido} \\
  \hline
  \caption[\sintesisResultadosCaption]{\sintesisResultadosCaption Fuente: Elaboración propia.}
  \label{tab:sintesis_resultados}
\end{longtable}
\endgroup
